\chapter{Exponential Decay of Correlations}\label{ch:correlations}

In the thermodynamic limit, spatial correlation functions of matrix product
states are governed by the subleading spectrum of the transfer map $\E_A$.
Subtracting the leading fixed-point contribution leaves a connected correlator
decaying exponentially (with polynomial factors $n^k$ for non-diagonal Jordan
blocks) at a rate defining the spatial correlation length. For tensors
generating parent Hamiltonians (Chapter~\ref{ch:parent_hamiltonian}), this
spectral gap controls ground-state correlations.

While Chapter~\ref{ch:spectral} analyzes the peripheral spectrum and overlap
decay, this chapter develops the one-point and two-point functionals, their
spectral decomposition, and quantitative decay bounds.

\section{Connected correlators from the transfer map}

\begin{definition}[Auxiliary-space one-point functional]\label{def:one_point_expectation}
    \lean{MPSTensor.onePointExpectation}
    \leanok
    \uses{def:transfer_map, thm:qpf}
    Fix an MPS tensor and a right fixed point of its transfer map.
    Write the tensor as $A$ and the fixed point as $\rho^R$.
    Under the hypotheses of Theorem~\ref{thm:qpf}, that fixed point exists
    and is unique up to scaling; we normalize it so that $\tr(\rho^R)=1$.
    For an auxiliary-space insertion $X\in\MN{D}$, define
    \begin{align}
        \langle X_0 \rangle & = \tr(X\rho^R). \notag
    \end{align}
\end{definition}

\begin{definition}[Auxiliary-space two-point functional at distance $n$]\label{def:two_point_expectation}
    \lean{MPSTensor.twoPointExpectation}
    \leanok
    \uses{def:transfer_map, def:one_point_expectation}
    For auxiliary-space insertions $X,Y\in\MN{D}$ and $n\ge0$, define
    \begin{align}
        \langle X_0 Y_n \rangle
         & = \tr\!\left(Y\,\E_A^n(X\rho^R)\right). \notag
    \end{align}
    Equivalently, insert $X$ at site $0$, propagate by $\E_A^n$, then insert
    $Y$ and take the trace.
    \begin{center}
        \tenkzkernel
        % Formula: \langle X_0Y_n\rangle
        %   = \tr\!\left(Y\,\E_A^n(X\rho^R)\right).
        % Source verdict: Cirac2021Matrix, Section II.B.3, supports the
        % correlator discussion but supplies no dedicated picture of this
        % contraction; the displayed definition is authoritative here.
        % For composite indices a=(k,l) and b=(i,j), the doubled map carries
        % (\E_A^n)_{ij,kl}; the east cup supplies (X\rho^R)_{kl}, and the
        % west cup contracts Y_{ji} with the output.  Thus the ink denotes
        % \sum_{i,j,k,l}Y_{ji}(\E_A^n)_{ij,kl}(X\rho^R)_{kl}.
        % The east cup is the applied-channel input; reading from the west
        % instead feeds the adjoint and gives the same trace pairing with Y.
        % Boundary signature: (virtual-west=0 | virtual-east=0 |
        % physical-up=0 | physical-down=0).
        \begin{tenkz}[rows={wire,wire}, cols=3,
                west={cup=$Y$}, east={cup=$X\rho^R$}]
            \tn[at=(1,2), skin=ring, wires=2]{\E_A^n}
        \end{tenkz}
    \end{center}
\end{definition}

\begin{definition}[Auxiliary-space connected correlator]\label{def:connected_correlator}
    \lean{MPSTensor.connectedCorrelator}
    \leanok
    \uses{def:two_point_expectation, def:one_point_expectation}
    The connected two-point function is
    \begin{align}
        C(X,Y;n)
         & = \langle X_0Y_n\rangle
        - \langle X_0\rangle\langle Y_0\rangle. \notag
    \end{align}
\end{definition}

\begin{definition}[Fixed-point-centered auxiliary insertion]%
    \label{def:correlator_traceless_part}
    \lean{MPSTensor.tracelessPart}
    \leanok
    For a matrix $\rho\in\MN{D}$ and an auxiliary insertion
    $X\in\MN{D}$, define
    \begin{align}
        Z_\rho(X)
         & = X\rho-\tr(X\rho)\rho.
        \label{eq:correlator_traceless_part}
    \end{align}
    When $\tr(\rho)=1$, this removes the trace component of $X\rho$.
\end{definition}

\begin{lemma}[Trace of the centered insertion]%
    \label{lem:trace_correlator_traceless_part}
    \lean{MPSTensor.trace_tracelessPart}
    \leanok
    \uses{def:correlator_traceless_part}
    If $\tr(\rho)=1$, then
    \begin{align}
        \tr\!\bigl(Z_\rho(X)\bigr) & =0.
    \end{align}
\end{lemma}

\begin{proof}
    \leanok
    \uses{def:correlator_traceless_part}
    By linearity of the trace,
    \begin{align}
        \tr\!\bigl(Z_\rho(X)\bigr)
         & =\tr(X\rho)-\tr(X\rho)\tr(\rho)
        =0.
    \end{align}
\end{proof}

\begin{theorem}[Fixed-point centering identity for connected correlators]%
    \label{thm:connected_correlator_traceless_reduction}
    \lean{MPSTensor.connectedCorrelator_eq_trace_transfer_tracelessPart}
    \leanok
    \uses{def:connected_correlator, def:correlator_traceless_part, def:transfer_map}
    Suppose that $\rho$ is a right fixed point of the transfer map,
    $\E_A(\rho)=\rho$.  Then, for every $n\geq0$,
    \begin{align}
        C(X,Y;n)
         & =\tr\!\left(Y\,\E_A^n\bigl(Z_\rho(X)\bigr)\right).
        \label{eq:connected_correlator_traceless_reduction}
    \end{align}
    When moreover $\tr(\rho)=1$, the preceding lemma shows that
    $Z_\rho(X)$ is traceless.  The identity isolates the fixed-point
    contribution before the remaining transfer iterates are analysed.
\end{theorem}

\begin{proof}
    \leanok
    \uses{def:connected_correlator, def:correlator_traceless_part}
    Since $\E_A(\rho)=\rho$, one has $\E_A^n(\rho)=\rho$ for every $n$.
    Linearity therefore gives
    \begin{align}
        \tr\!\left(Y\,\E_A^n\bigl(Z_\rho(X)\bigr)\right)
         & =
        \tr\!\left(Y\,\E_A^n(X\rho)\right)
        -\tr(X\rho)\tr(Y\rho),
    \end{align}
    which is the definition of $C(X,Y;n)$.
\end{proof}

\section{Spectral expansion and exponential decay}

\begin{theorem}[Sum of exponentials (spectral expansion)]%
    \label{thm:correlator_sum_exponentials}
    \lean{MPSTensor.connectedCorrelator_eq_sum}
    \leanok
    \uses{def:connected_correlator}
    If coefficients $c_j\in\C$ and eigenvalues $\lambda_j\in\C$, for
    $j=1,\ldots,D^2-1$, satisfy, for all $n\ge 0$,
    \begin{align}
        C(X,Y;n)
         & = \sum_{j=1}^{D^2-1} c_j\,\lambda_j^n,
        \label{eq:correlations_spectral_expansion}
    \end{align}
    then the connected correlator equals the stated sum of exponentials.
    When the subleading spectrum of the transfer map is diagonalizable, the
    connected correlator expands as a sum of pure exponentials
    $C(X,Y;n)=\sum_{j=1}^{D^2-1} c_j\,\lambda_j^n$
    \cite[Section~II.B.3]{Cirac2021Matrix}; Jordan blocks can produce polynomial
    factors $n^k\lambda^n$ when the transfer map is defective.
\end{theorem}

\begin{proof}
    \leanok
    Given the expansion hypothesis
    (\ref{eq:correlations_spectral_expansion}), the equality is immediate.
    Under the corresponding diagonalizability hypothesis, the expansion
    follows from the spectral decomposition of $\E_A^n$ restricted to the
    complement of the fixed-point eigenspace. The leading eigenvalue
    $\lambda_1=1$ contributes
    $\langle X_0\rangle\langle Y_0\rangle$, which is subtracted in the
    connected correlator, leaving the sum over subleading eigenvalues.
    This is the spectral mechanism described in
    \cite[Section~II.B.3]{Cirac2021Matrix}.
    % Constructing the eigendecomposition coefficients from $\E_A$
    % is tracked in Issue~\#1447.
\end{proof}

\begin{theorem}[Exponential decay bound]%
    \label{thm:correlator_exponential_bound}
    \lean{MPSTensor.connectedCorrelator_bound}
    \leanok
    \uses{def:connected_correlator}
    If a constant $C_{XY}\in\R$ and $\lambda_2\in\C$ satisfy, for all
    $n\ge0$,
    \begin{align}
        |C(X,Y;n)|
         & \le C_{XY}\,|\lambda_2|^n,
        \label{eq:correlations_decay_bound}
    \end{align}
    then the connected correlator satisfies that exponential decay bound.
    Combining the spectral expansion with the triangle inequality yields this
    bound with $|\lambda_2|<1$ whenever the transfer map is primitive or
    injective \cite[Section~II.B.3]{Cirac2021Matrix}.
\end{theorem}

\begin{proof}
    \leanok
    Given the bound hypothesis (\ref{eq:correlations_decay_bound}), the
    estimate is immediate.
    For a pure-exponential expansion, the source
    \cite[Section~II.B.3]{Cirac2021Matrix} obtains such a bound by applying
    the triangle inequality and using
    $|\lambda_j|\le|\lambda_2|$ for all subleading eigenvalues; when the
    transfer-map gap gives $|\lambda_2|<1$, this is a decaying bound.
    % Deriving the bound from the eigendecomposition of $\E_A$
    % is tracked in Issue~\#1447.
\end{proof}

\begin{remark}[Transfer-map spectral hypothesis]%
    \label{rem:spectral_hypothesis}
    \uses{thm:correlator_exponential_bound, thm:transfer_operator_gap_irr_tp, thm:transfer_map_gap_wielandt}
    The condition $|\lambda_2|<1$ turns the bound in
    Theorem~\ref{thm:correlator_exponential_bound} into exponential decay. It
    is a transfer-map gap condition: all non-leading eigenvalues of $\E_A$ lie
    strictly inside the unit disk. It is distinct from the Hamiltonian spectral gap of
    Chapter~\ref{ch:parent_hamiltonian_commuting_and_gap}.
    The complementary transfer-map gap $\rho(\E_A-P)<1$ is established for
    primitive channels
    (the corresponding theorem in the companion quantum-channel volume~\cite[``Complementary transfer-map gap for primitive maps'']{QICLean2026Blueprint}).
    The mixed-transfer gap needed for separation of non-equivalent canonical
    blocks is supplied for irreducible trace-preserving tensors
    (Theorem~\ref{thm:transfer_operator_gap_irr_tp}), and quantitative
    single-tensor transfer-map bounds for injective tensors are recorded in
    Theorem~\ref{thm:transfer_map_gap_wielandt}.

    These transfer-map gap results give exponential convergence of the
    complementary transfer map at every rate strictly exceeding its
    spectral radius.
\end{remark}

\begin{definition}[Correlation length]\label{def:correlation_length}
    \lean{MPSTensor.correlationLength}
    \leanok
    \uses{def:transfer_map}
    For every $\lambda_2\in\C$, define the associated correlation-length
    quantity by
    \begin{align}
        \xi
         & = -\frac{1}{\log|\lambda_2|}. \notag
    \end{align}
    For a nonzero subleading eigenvalue in the physical range
    $0<|\lambda_2|<1$, this quantity is positive. If all subleading spectral
    values vanish and correlations therefore vanish after finitely many
    transfer steps, the total definition gives the limiting value $\xi=0$.
\end{definition}

\begin{lemma}[Correlation length is positive]\label{lem:correlation_length_pos}
    \lean{MPSTensor.correlationLength_pos}
    \leanok
    \uses{def:correlation_length}
    If $|\lambda_2| < 1$ and $\lambda_2 \ne 0$, then $\xi > 0$.
\end{lemma}

\begin{proof}
    \leanok
    Since $0<|\lambda_2|<1$, we have $\log|\lambda_2|<0$, so
    $\xi=-1/\log|\lambda_2|>0$.
\end{proof}

\section{Quantitative transfer-map gap bounds}

The qualitative bound $\rho(\E_A-P)<1$ for primitive channels
\cite[``Complementary transfer-map gap for primitive maps'']{QICLean2026Blueprint}
admits quantitative refinements in terms of geometric constants and decay rates
for the transfer map $\E_A$.

\begin{theorem}[Exponential convergence of injective primitive channels]%
    \label{thm:exponential_convergence_primitive}
    \lean{MPSTensor.exponential_convergence_of_primitive}
    \leanok
    \uses{thm:primitive_overlap}
    Fix an injective, trace-preserving MPS tensor $A$ with positive
    definite fixed point $\rho$.
    Let
    \begin{align}
        P(X)
         & = \frac{\tr(X)}{\tr(\rho)}\rho
        \label{eq:correlations_fixed_point_projection}
    \end{align}
    be the corresponding fixed-point projection.
    Then there exist $C > 0$ and $0 < \delta \le 1$ such that, for all
    $n\ge 0$ and $X\in\MN{D}$,
    \begin{align}
        \|\E_A^n(X)-P(X)\|
         & \le C(1-\delta)^n\|X\|.
        \label{eq:correlations_channel_convergence}
    \end{align}
    The transfer-map gap $\delta$ exists by primitivity
    (the corresponding theorem in the companion quantum-channel volume~\cite[``Complementary transfer-map gap for primitive maps'']{QICLean2026Blueprint}, the complementary transfer-map
    gap).
\end{theorem}

\begin{proof}
    \leanok
    Injectivity implies primitivity of the transfer map.
    The complementary spectral-radius bound $\rho(\E_A-P)<1$ then follows
    from the primitive transfer-map gap theorem. Choose
    $r$ with $\rho(\E_A-P)<r<1$. The Gelfand formula yields
    $\|(\E_A-P)^n\|\le Cr^n$ for a suitable constant $C>0$.
    Since $P$ is the fixed-point projection, $\E_A P=P\E_A=P$, and hence
    $\E_A^n-P=(\E_A-P)^n$ for $n\ge 1$. The case $n=0$ is absorbed by
    increasing $C$ if necessary. Setting $\delta=1-r$ gives
    (\ref{eq:correlations_channel_convergence}).
\end{proof}

\begin{theorem}[Correlation length bound]\label{thm:correlation_length_bound}
    \lean{MPSTensor.correlation_length_bound}
    \leanok
    \uses{def:correlation_length, thm:exponential_convergence_primitive, def:injective}
    For an injective trace-preserving MPS tensor, there exist
    $C>0$ and $\xi>0$ such that, for all $n\ge 0$ and all
    traceless $X\in\MN{D}$, i.e.\ $\tr(X)=0$,
    \begin{align}
        \|\E_A^n(X)\|
         & \le C e^{-n/\xi}\|X\|.
        \label{eq:correlations_traceless_decay}
    \end{align}
    Traceless matrices lie in $\ker P$, where $P$ is the fixed-point
    projection (\ref{eq:correlations_fixed_point_projection}).
    Since $\E_A-P$ has spectral radius strictly less than~$1$, choose any
    $r$ with $\rho(\E_A-P)<r<1$. The Gelfand formula then gives exponential
    decay with $\xi=-1/\log r$. Choosing a rate strictly above the spectral
    radius is necessary in general because a Jordan block at the spectral
    radius can prevent a uniform bound proportional to
    $\rho(\E_A-P)^n$.
\end{theorem}

\begin{proof}
    \leanok
    For traceless $X$, we have $P(X)=0$, so
    $\E_A^n(X)=(\E_A-P)^n(X)$.
    Applying the Gelfand-formula estimate with
    $\rho(\E_A-P)<r<1$ gives $\|\E_A^n(X)\|\le Cr^n\|X\|$.
    With $\xi=-1/\log r>0$, the identity $r^n=e^{-n/\xi}$ yields
    (\ref{eq:correlations_traceless_decay}).
\end{proof}

\begin{theorem}[Transfer-map gap from injectivity]%
    \label{thm:transfer_map_gap_wielandt}
    \lean{MPSTensor.transfer_map_gap_of_injective}
    \leanok
    \uses{thm:transfer_operator_gap}
    For an injective trace-preserving MPS tensor, there exists
    $\delta>0$ such that every eigenvalue $\mu\ne 1$ of the
    transfer map satisfies $|\mu|\le 1-\delta$.
\end{theorem}

\begin{proof}
    \leanok
    Injectivity implies eventual full Kraus rank. This yields primitivity,
    hence a complementary transfer-map gap. Since the transfer map has
    only finitely many eigenvalues, one may choose a uniform
    $\delta>0$ separating $1$ from the remaining spectrum.
\end{proof}

\section{Relation to parent Hamiltonians}

\begin{remark}\label{rem:correlations_parent_hamiltonian}
    \uses{def:parent_hamiltonian, lem:parent_hamiltonian_annihilates, def:correlation_length, thm:correlator_exponential_bound}
    Suppose that the MPS tensor $A$ generates a parent Hamiltonian $H_N(A,L)$
    (Definition~\ref{def:parent_hamiltonian}) whose ground state is the matrix
    product vector (Lemma~\ref{lem:parent_hamiltonian_annihilates}), and that
    physical observable insertions reduce to the auxiliary functionals defined
    above. If the hypotheses of
    Theorem~\ref{thm:correlator_exponential_bound} hold, then the corresponding
    ground-state correlations decay with finite correlation length $\xi$.
\end{remark}
